\chapter{Einleitung}
 \section{Motivation}
Sensorarray beschreiben

\begin{figure}[ht!]
 \centering
 \includegraphics[width=0.8\textwidth]{img/AblaufFourier.pdf}
 \caption{Ablauf der Signalvorverarbeitung~\autocite[9]{krrts2017freqfilt}}
 \label{pic:AblaufFourier}
\end{figure}


\section{Stand der Technik}
Der verwendete Prozess ist mit $\SI{350}{nm}$ im Vergleich zu modernen Prozessen mit $\SI{20}{\nm}$ oder weniger Strukturbreite um die Größenordnung 20 größer. Entsprechend handelt es 
sich um einen relativ alten Prozess.

Kurze Beschreibung zu Standardzellen.


\section{Ziel dieser Arbeit}
Im Rahmen des \gls{isar}-Projekts der HAW Hamburg soll zur Signalvorverarbeitung einer Matrix von Magnetsensoren  eine 
%\gls{2d-dft} 
zweidimensionale diskrete Fouriertransformation (2D-DFT) in VHDL implementiert werden. Mit der 
%\gls{2d-dft} 
2D-DFT sollen relevante Signalanteile identifiziert werden, um so den Informationsgehalt der Sensorsignale auf relevante Anteile zu reduzieren. Die Sensoren basieren auf 
dem \gls{amr2} bzw. in einem späteren Schritt \gls{tmr2}.

In einem Text zitiert dann so~\autocite[10-20]{krey2015systemarchitektur} und blabla.

Winkelberechnung als Ziel beschreiben. und noch ein Zitat~\autocite[8-9]{krrts2017freqfilt} noch eins~\autocite[1-2]{sabotta2013Masterthesis}

